% ===================================================================
% CHAPTER 8: LIMITATIONS
% ===================================================================
\chapter{Ограничувања}
\label{chap:limitations}

Ова поглавје ги изнесува ограничувањата на истражувањето. Тоа е напишано намерно исцрпно и
намерно против сопствениот предмет: секоја слабост што ревизијата на податоците ја откри е
наведена заедно со проверката што покажува колку таа влијае врз заклучокот. Наодите што
\textbf{не} преживеале проверка се повлечени, а не ублажени.

\section{Ограничувања на статистичката моќ}

Дизајнот е спарен по почетна вредност, со $N = 10$ пара. При тој обем, двостраниот
Вилкоксонов тест на рангирани знаци има \textbf{долна граница} од $p = 0{,}0020$. Тоа значи две
работи истовремено: сите четири примарни показатели ја достигнуваат максималната достижна
значајност, но и дека $p$-вредноста сама по себе носи малку информација --- таа не може да
се разликува меѓу „голем ефект“ и „огромен ефект“.

Затоа тврдењата ги носат \textbf{големината на ефектот и интервалот на доверба}. Cliff's
$\delta = \pm 1{,}00$ означува целосно раздвојување: секое извршување со едниот пристап е
подобро од секое извршување со другиот. Тоа е посилно тврдење од која било $p$-вредност при
овој обем, бидејќи не зависи од претпоставки за распределбата. Целосните податоци по
почетна вредност се објавени, за читателот да може да ја процени варијабилноста самостојно.

\section{Спарувања извршени во различни сесии}

Ова е најсериозната закана врз валидноста и затоа е разгледана прва.

Од десетте пара, \textbf{шест} се извршени во иста сесија (почетни вредности 42, 46, 48, 49, 50,
51), а \textbf{четири} во различни сесии, оддалечени до околу шест дена (почетни вредности 43,
44, 45, 47). Кај тие четири пара, амбиенталните услови се разликувале \emph{во рамките на самиот
пар}, што е токму она што спарениот дизајн треба да го спречи.

Проверката на робусноста го решава прашањето. Ако кампањата се подели на двете подмножества,
ефектот останува непроменет:

\begin{table}[htbp]
\centering
\caption{Робусност на времето до обновување во однос на спарувањето по сесија.}
\label{tab:robustness-session}
\begin{tabular}{lccccc}
\toprule
Подмножество & Наивно (мед.) & SISA (мед.) & Медијана на односот & $\delta$ & $p$ \\
\midrule
Иста сесија ($n = 6$) & 251,6 s & 31,3 s & 7,9$\times$ & $\pm 1{,}00$ & 0,031 \\
Различна сесија ($n = 4$) & 240,5 s & 27,3 s & 8,8$\times$ & $\pm 1{,}00$ & 0,125 \\
\bottomrule
\end{tabular}
\end{table}

Во \textbf{двете} подмножества раздвојувањето е целосно и нема ниту едно преклопување.
$p$-вредностите се граничните вредности при $n = 6$ и $n = 4$ соодветно, не показатели на
послаб ефект. Освен тоа, вистинската контрола врз амбиенталните услови не е спарувањето по
сесија, туку \textbf{задолжителното ладење}: секое мерно извршување започнува од стабилизирано
термичко плато. Разликата од околу 210 s меѓу краковите е за цел ред на големина поголема
од кое било термичко колебање од неколку десетици секунди.

Заклучокот е дека прекуссесиското спарување не може да го преврти резултатот. Сепак,
преправањето на четирите пара во иста сесија останува препорачливо подобрување, и се
наведува како такво, а не се прикажува како непотребно.

\section{Амбиенталната температура не е снимена}

Протоколот предвидуваше амбиенталната температура да се запишува преку сензор DS18B20, со
температурата на платото како резервен показател запишан рачно.

\textbf{Ниту едно од двете не е спроведено.} Сите дваесет телеметриски датотеки ја имаат
десетколонската заглавна линија, без колоната \texttt{Ambient\_C} --- сензорот бил додаден во
телеметриската скрипта по извршувањата. Ниту едно извршување нема придружна датотека со
белешки. Амбиенталната температура затоа е \textbf{незапишана за целата кампања}, која се
протега од крајот на јули до почетокот на август 2026.

Последицата е точно определена: \textbf{апсолутните} термички вредности (врвна температура,
секунди со придушување) не се споредливи меѓу извршувања изведени во различни денови, и не
треба да се толкуваат како својство единствено на работното оптоварување. \textbf{Спарените
разлики} остануваат валидни, бидејќи задолжителното ладење ја изедначува почетната
состојба, а во шест од десетте пара двата крака се извршени во иста сесија. Тврдењата во
\hyperref[chap:results]{Дел~\ref*{chap:results}} се формулирани врз спарените разлики, не врз апсолутните вредности, токму
поради ова ограничување.

\section{Резидуален поднапон}

Протоколот бара регистарот за придушување да чита \texttt{0x0} пред секое извршување, по
замената со кабел од 5 A. Проверката на телеметријата покажува дека тоа \textbf{не е целосно
постигнато}: живиот бит за поднапон се појавува во околу половина од извршувањата.

Мерењето по извршување го определува обемот прецизно. Низ целото извршување поднапонот
достигнува до 45 s (наивно, почетна вредност 49), но во \textbf{прозорецот на обновувањето} --- единствениот што се мери --- вредностите се меѓу 0 и 9 s. Најлошиот случај е повторно
почетната вредност 49, со 9 s од вкупно 237,6 s, односно околу 3,8\,\% од прозорецот. Кај
шеснаесет од дваесетте извршувања поднапонот во Фаза 3 е нула или една секунда. Појавата е
концентрирана во Фазите 1 и 4, каде оптоварувањето е долготрајно.

Две причини зошто тоа не го загрозува резултатот:

\begin{enumerate}
  \item \textbf{Не го надувува показателот за придушување.} Поднапонот е битот 0 во регистарот, а мерката за термичко придушување ги брои битовите 1–3. Двете се читаат одвоено.
  \item \textbf{Не ги предизвикува падовите на тактот.} Кај почетните вредности 42 и 45 наивниот крак паднал на 1500 MHz со \textbf{нула} секунди поднапон. Падовите на фреквенцијата се затоа термички, не последица на напојувањето.
\end{enumerate}

Ограничувањето се известува како квантифицирано и маргинално, а тврдењето во протоколот
дека регистарот „мора да биде 0x0“ се коригира во „поднапонот е сведен на маргинална појава
надвор од мерниот прозорец“.

\section{Што точно мери показателот за придушување}

Секундите со придушување се бројат врз отчитите означени со маркерот на Фаза 3. Тој
прозорец го опфаќа \textbf{целото повикување} на обновувањето --- подигањето на контејнерот,
пресметката и расчистувањето --- и затоа е неколку секунди подолг од чистото пресметковно
време што го известува модулот за обновување.

Последицата е видлива кај почетните вредности 50 и 51, каде SISA бележи 49 и 47 секунди
придушување наспроти времиња на обновување од 43,8 и 43,1 s. Тоа не е недоследност, туку
дефиниција: показателот е „секунди со придушување во прозорецот на обновувањето“, не
„секунди со придушување при пресметката“. Бидејќи истата дефиниција важи за двата крака,
спарената разлика не е пристрасна; апсолутните вредности, пак, не смеат да се толкуваат
како чисто пресметковно време.

\section{Детерминизам на пресметката, но не и на часовникот}

Протоколот тврди дека извршувањата се детерминистички. Тоа тврдење важи за
\textbf{пресметката}, не за \textbf{мерењето на времето}. Повторено извршување на обновувањето дава
идентична загуба до последната децимала (на пример, 0,014226 во двата обида кај почетната
вредност 51), но времетраењето се разликува (263,4 наспроти 261,5 s), бидејќи времето е
физичко мерење осетливо на моменталната термичка состојба и на тактот.

Ова е очекувано и коректно, но треба да се наведе експлицитно, бидејќи „детерминистички
експеримент“ инаку би значело дека и времињата треба да се повторат точно.

\section{Ненаправени помошни извршувања}

Три помошни серии беа предвидени во претходната регистрација и \textbf{не се спроведени}:

\begin{itemize}
  \item \textbf{Чисти референтни извршувања} (\texttt{POISON\_MODE=off}). Без нив, хипотезата H5 во нејзината првично регистрирана форма --- „обновениот модел е во рамките на 2 процентни поени од модел обучен врз чисти податоци“ --- \textbf{не може да се оцени како што беше специфицирана}. Затоа H5 во \hyperref[chap:results]{Дел~\ref*{chap:results}} е известена во изменета форма: како споредба меѓу двата крака и наспроти тривијалниот класификатор на мнозинската класа, а не наспроти чиста горна граница. Оваа измена во однос на претходната регистрација се наведува отворено.
  \item \textbf{Студија на чувствителност со скалиран модел.} Ова беше предвидениот начин да се лоцира големината на моделот при која влезно-излезните операции стануваат доминантни. Без неа, наодот за H4 останува ограничен на еден единствен ред на големина на моделот.
  \item \textbf{Извршување под WAN-услови.} Предвидено да покаже дека резултатот се задржува при реалистично мрежно доцнење. Мерниот прозорец на Фаза 3 е чисто локална пресметка на јазолот и не вклучува мрежна комуникација, па не се очекува влијание --- но тоа е \textbf{очекување, не измерен факт}.
\end{itemize}

\section{Ограничувања на проверката на заборавањето}

Егзактноста на заборавањето овде почива на \textbf{конструкцијата}: отстранетите редови
никогаш не влегуваат во повторно обработените составни модели, ниту во моделот преобучен од
нула. Автоматскиот тест го потврдува детерминизмот на таа конструкција преку битова
споредба на обновениот модел меѓу две извршувања, и потврдува дека точката на враќање
наназад е избрана пред компромитирањето.

Емпириската проверка, меѓутоа, е \textbf{незаклучна}. Пробата за заклучување членство ја
споредува загубата врз отстранетите примери со онаа врз примери што никогаш не биле
употребени. Со модел од 12\,865 параметри обучен врз 800\,000 реда, ниту контролната група на
задржани примери не покажува мерлив сигнал за членство --- што значи дека пробата нема
разделна моќ, а не дека моделот заборавил. Ова се известува како незаклучен резултат.
Емпириска потврда на заборавањето би барала или значително поголем модел, или
поосетлива метода за напад.

\section{Ограничувања на генерализацијата}

\begin{itemize}
  \item \textbf{Еден уред, една картичка.} Сите мерења се изведени на еден примерок Raspberry Pi 5 со една SD-картичка. Термичкото однесување зависи од конкретниот примерок, од неговото куќиште и од распоредот на воздухот; влезно-излезните карактеристики зависат од конкретната картичка и од степенот на нејзината искористеност. Трудот треба да се чита како \textbf{студија на случај врз еден уред}, не како карактеризација на класата уреди.
  \item \textbf{Еден модел, една задача.} Заклучоците важат за повеќеслоен перцептрон со 12\,865 параметри што решава бинарна класификација. Особено наодот за H4 е директна последица на малата големина на моделот.
  \item \textbf{Едно податочно множество.} Изразената класна нерамнотежа од околу 96,5\,\% напади е својство на TON\_IoT и на начинот на кој е издвоено тест-множеството. Наодот за колапсот во мнозинската класа најверојатно зависи од таа нерамнотежа и не смее да се пренесе на избалансирани задачи без повторно мерење.
  \item \textbf{Откривањето е надвор од опсегот.} Компромитираните примери се дадени преку оракул-маска. Во реален систем идентификацијата е сама по себе тежок проблем и нејзината цена не е вклучена во ниту еден од измерените трошоци.
\end{itemize}

\section{Ограничувања на самата конструкција}

\begin{itemize}
  \item \textbf{Приближно глобално заборавање.} Двата пристапа даваат егзактно \textbf{локално} заборавање. Придонесот што е веќе агрегиран во глобалниот модел не може да се отстрани од страната на клиентот при FedAvg без состојба. Заклучокот важи за таа поставка; системи каде серверот чува историја на ажурирања имаат други можности и друга цена.
  \item \textbf{Две отстапувања од класичниот SISA.} Составните модели никогаш не ги преземаат глобалните тежини --- што е условот гаранцијата при враќање наназад да важи --- а FL-ажурирањето е параметарска средна вредност наместо ансамбл од предвидувања. Второто отстапување е \textbf{директната причина} за колапсот во мнозинската класа утврден во \hyperref[chap:results]{Дел~\ref*{chap:results}}, §7.5. Тоа значи дека измерената корисност е својство на \textbf{оваа интеграција} на SISA во FedAvg, не на SISA воопшто.
  \item \textbf{Предноста на SISA зависи од локалноста на компромитирањето.} Целиот добиток произлегува од тоа што труењето е ограничено на доцни сегменти на еден шард; при компромитирање во ран сегмент предноста исчезнува. Ова е измерено и разгледано како услов на важење во \hyperref[chap:results]{Дел~\ref*{chap:results}}, §7.4, па овде се наведува само како потсетник дека резултатот е врзан за моделот на закана, а не универзален.
  \item \textbf{Мала непрецизност во тежинското вреднување.} При повторното приклучување, клиентот со SISA го пријавува целиот регион за обука (800\,000 реда) како број примери кон FedAvg, иако околу 63\,000 редови се исклучени. Придонесот на јазолот затоа е малку преценет во агрегацијата во Фаза 4 --- ефект од околу 8\,\% врз тежината на еден од четирите клиенти.
\end{itemize}

\section{Методолошки ситници што треба да се наведат}

\begin{itemize}
  \item \textbf{Стандардизацијата е применета врз целата партиција}, вклучувајќи ги и редовите што подоцна стануваат валидациски и дел од глобалното тест-множество. Не постои посебен стандардизатор изведен само врз податоците за обука. Тоа е благ пренос на информација, идентичен за двата крака, па не влијае на споредбата, но апсолутните вредности на корисноста се благо оптимистички.
  \item \textbf{Енергијата во Фаза 1 недостасува за една почетна вредност} (45), каде дневникот на потрошувачката бил рестартиран и фазата 1 паднала надвор од покриениот интервал. Таа величина е описна и не влегува во спарената статистика; енергијата во мерниот прозорец на Фаза 3 е целосна за сите дваесет извршувања.
  \item \textbf{Мрежното доцнење во примарните извршувања е занемарливо}, бидејќи целата федерација работи преку локална мрежа. Тоа е реалистично за фабрички или домашен пораб, но не и за географски распределени федерации.
\end{itemize}

\section{Уделот на труењето не е константен низ кампањата}
\label{subsec:poison-fraction}

Претходно регистрираниот дизајн го замрзнува уделот на превртени напади на \textbf{1,0}. Ревизијата
на манифестите покажува дека тоа важи за \textbf{осум} од десетте почетни вредности, но не за сите:

\begin{table}[htbp]
\centering
\caption{Големина на компромитираното множество по почетна вредност.}
\label{tab:poison-fraction}
\begin{tabular}{lll}
\toprule
Почетни вредности & Компромитирани примери & Импликуван удел \\
\midrule
42--49 & 63\,060--63\,155 (варијација 0,15\,\%) & 1,0 \\
\textbf{50, 51} & \textbf{31\,517 и 31\,561} & \textbf{0,5} \\
\bottomrule
\end{tabular}
\end{table}

Односот на двете групи изнесува 2,001, што при варијација од само 0,15\,\% во рамките на првата
група не остава простор за друго објаснување освен различна вредност на параметарот
\texttt{--poison-fraction}. Најверојатното потекло е подразбираната вредност \texttt{0.5} во
\texttt{experiments/prepare\_edge\_data.py}, која никогаш не беше усогласена со замрзнатиот протокол и
која се применува кога знаменцето не е зададено при подготовката на податоците.

\textbf{Што ова значи, и што не значи.} Во рамките на еден пар, двата крака работат врз \emph{истото}
компромитирано множество --- подготовката на податоците се врши еднаш по почетна вредност и
двата крака ја користат истата датотека. Спарената разлика затоа останува валидна за сите
десет пара, а тоа е величината на која се засноваат сите тврдења во \hyperref[chap:results]{Дел~\ref*{chap:results}}. Она што
е загрозено е \textbf{хомогеноста на третманот} низ кампањата: две од десетте почетни вредности
отстрануваат половина помало множество, и затоа на двата крака им остануваат повеќе редови за
обработка. Тоа е и измерената причина за нивните повисоки апсолутни времиња (§7.4), и токму
затоа тие вредности се изнесени со објаснување наместо да бидат припишани на амбиенталните
услови.

Проверката е тривијална и е додадена во \texttt{analysis/analyze\_runs.py}, кој сега предупредува кога
големината на компромитираното множество не е константна низ почетните вредности. Исправката
во кодот --- подразбираната вредност да се усогласи со протоколот --- е наведена во
\texttt{07-глава05-методологија-исправки.md}, И1. Повторувањето на почетните вредности 50 и 51 со
удел 1,0 би ја отстранило оваа забелешка целосно и е препорачливо, но не е неопходно за
важноста на спарените резултати.

\section{Оптималниот праг е оракул, не проценка за распоредување}

Наодот во \hyperref[chap:results]{Дел~\ref*{chap:results}}, §7.5 --- дека со подесен праг SISA достигнува избалансирана точност
0,709 наспроти 0,671 за наивниот пристап --- има ограничување што мора да се наведе јасно,
бидејќи инаку бројката се чита како ветување.

Прагот е избран така што ја \textbf{максимизира избалансираната точност врз истото множество врз
кое потоа се известува резултатот}. Тоа е оракул: тој ја дава \textbf{горната граница} на она што
калибрацијата може да го постигне, а не перформансата што би се добила во распоредување.
Реален систем го избира прагот однапред, врз податоци што не ги гледал при оценувањето, па
вистинската вредност е нужно пониска --- колку пониска, овој труд не мери.

Соодветната постапка е прагот да се бира на издвоено валидациско множество и да се оценува на
дисјунктно тест-множество. Тоа е директно изводливо од постојните артефакти и е наведено меѓу
идните чекори (\hyperref[chap:future]{Дел~\ref*{chap:future}}, §10.1).

Второ ограничување од истата анализа: сите дваесет модела се оценети врз \textbf{едно заедничко}
тест-множество (она подготвено со почетната вредност 51), бидејќи тест-множеството се
подготвува по почетна вредност и се презапишува под исто име. Тоа множество е дисјунктно од
множеството за обука само на \textbf{едната} почетна вредност; кај останатите девет дел од
редовите можеби учествувале во нивната обука. Спарената споредба не е засегната --- двата крака
на едно семе се оценети со идентичен критериум --- но апстолутните вредности на корисноста се
благо оптимистички. \texttt{analysis/evaluate\_utility.py} сега предупредува за ова и поддржува
оценување со тест-множество по семе преку \texttt{--test-template}; спроведувањето бара повторно
генерирање на кешовите за сите десет почетни вредности.

